% !PRINTED 706
\weberpage{706}
Since this function can accordingly contain no negative powers of $z$, the theorem to be proved follows from (13):
\[
\overset{i}{\Sigma} a^{(i)}_{-1} = 0,
\]
which no longer has anything to do with the particular assumption about the variable $z$.

As a special application it follows from this that a differential of the second kind has no residue, and that the two residues of the differentials of the third kind are equal and opposite.

\gesperrt{A proper differential} $d\sigma$, as we have already seen, \gesperrt{has no residue}, and the residues of the \glqq\gesperrt{logarithmic differential}\grqq{} $d\sigma:\sigma$ are integers, namely the order numbers in $z$.

We will conclude these considerations with the proof of the following theorem:

4. \gesperrt{If we call differentials of the second kind linearly dependent when a linear combination of them and of differentials of the first kind with constant coefficients is equal to a proper differential, then at most $p+1$ differentials of the second kind are always linearly dependent.}

To prove this, we take a $p$-gon of the second kind:
\weq{18}{\mathfrak P_1 \mathfrak P_2 \ldots \mathfrak P_p}
with $p$ distinct points. If then $\mathfrak P$ is a given point, we set:
\weq{19}{\mathfrak U = \mathfrak P \mathfrak P_1 \mathfrak P_2 \ldots \mathfrak P_r}
% [sic] printed P_r (Fraktur) as the last factor of (19), a misprint for P_p: the p-gon (18) is P_1 P_2 ... P_p
and obtain a class $B$ whose divisor $\mathfrak P$ in any case does not divide, since otherwise $\mathfrak P_1 \mathfrak P_2 \ldots \mathfrak P_n$
% [sic] printed P_1 P_2 ... P_n (Fraktur), a misprint for P_1 P_2 ... P_p: the polygon meant is the p-gon (18)
would have to be of the first kind by $\S\,199$, 7. We can therefore determine a function in $\Omega$:
\weq{20}{z = \frac{\mathfrak A}{\mathfrak U}}
in whose lower polygon the point $\mathfrak P$ is in any case contained. We now form the proper differential
\weq{21}{dz = \frac{\mathfrak Z}{\mathfrak U^2},}
whose numerator $\mathfrak Z$ (the branch polygon with respect to $z$) the point $\mathfrak P$ does not divide, because it divides $\mathfrak U$ only simply, and if we express $dz$, according to $\S\,200$, 4., by differentials of the first, second and third
%CONT>>
